\section{When a history-uniform task mark exists}
\label{sec:fiber}
\begin{theorem}[Conditional minimax radius criterion]\label{thm:fiber}
Under the fiber regularity in Section~\ref{sec:setup}, \eqref{eq:canonical-mark} defines bounded Borel functions with a unique center. The following are equivalent, up to the displayed constants.

\smallskip\noindent
(i) There are Borel $f,w\ge0$ and $0<a\le A<\infty$ such that, for every $(\mu,s)\in\mathcal F_x$ and $r\in H$,
\[
 a\{w+\|f-r\|^2\}\le s+\|\mu-r\|^2
       \le A\{w+\|f-r\|^2\}.
\]
\noindent
(ii) $w_*(x)\le K\underline s(x)$ uniformly for some finite $K$.

Condition (i) implies (ii) with $K=A/a$. Condition (ii) implies (i) for the canonical mark with $a=(3K)^{-1}$, $A=2$. Zero-radius fibers are included without division by zero. This is a criterion for the stated all-prediction, all-history quadratic comparison, not a necessary condition for every finite-resource risk theorem.
\end{theorem}
\begin{proof}
Each function $r\mapsto s+\|\mu-r\|^2$ is strongly convex with the same quadratic part. Hence
\[
 g_x(\lambda r+(1-\lambda)r')
 \le\lambda g_x(r)+(1-\lambda)g_x(r')
       -\lambda(1-\lambda)\|r-r'\|^2.
\]
It is uniformly Lipschitz on compact $H$, uniformly over $x$. Existence follows by compactness and uniqueness by this inequality. To see Borel measurability without an implicit selection assumption, minimize $g_x$ over successive fixed finite nets of a countable dense subset of $H$, breaking ties by index. These minimizers are Borel. Their objective errors tend uniformly to zero by Lipschitz continuity, and strong convexity on the segment to the unique minimizer implies convergence to it. Therefore $f_*$ and $w_*$ are Borel.

Under (i), putting $r=f$ gives $g_x(f)\le A w$. Putting $r=\mu$, separately for every pair in the fiber, gives $aw\le s$. Thus
\[
 w_*\le g_x(f)\le Aw\le(A/a)\underline s.
\]
The quantifier over every fixed $r\in H$ permits this pointwise substitution; it does not give the observer access to $\mu$.

Conversely let (ii) hold. For every conditional pair,
\[
 s+\|\mu-f_*\|^2\le w_*,\qquad s\ge w_*/K.
\]
The upper comparison follows from
\[
 s+\|\mu-r\|^2
 \le2\{s+\|\mu-f_*\|^2\}+2\|f_*-r\|^2
 \le2\{w_*+\|f_*-r\|^2\}.
\]
For the lower,
\[
 w_*+\|f_*-r\|^2
 \le3w_*+2\|\mu-r\|^2
 \le3K\{s+\|\mu-r\|^2\}.
\]
If $w_*=0$, all pairs have $s=0$ and $\mu=f_*$, and the comparison is immediate. Also $w_*\ge\underline s$, so $K\ge1$ can always be used on positive fibers.
\end{proof}

\subsection*{What the criterion distinguishes}
An exact task identity has a singleton fiber $(f,w)$ and directly yields the mark. A primitive minorization or positivity-improving channel can instead bound all conditional variances below and conditional means above, producing a sufficient certificate. A calibration supplied by an external protocol must pay for that protocol when it forms part of the experiment. None of these cases follows from positivity of an arbitrary normalized kernel alone.

\begin{example}[Repeated information can destroy a current-report mark]\label{ex:calibration-failure}
Let $E\in\{-1,1\}$ be fixed, initially fair, and let independent reports satisfy $\Pp(V=E)=1-\gamma>1/2$. Set the audit to $Z=E$. Repeated positive reports drive $\Pp(E=-1\mid V_1=\cdots=V_n=1)$ to zero. Histories ending in the same current report $V_n=1$ also include histories with a substantial negative posterior. Their fiber has least variance zero but positive minimax radius. Thus no positive uniform quadratic current-report mark exists. Enlarging the predictive state to retain the posterior is a different, separately charged implementation; the failed mark cannot be rescued by declaring its formal support acquired.
\end{example}

\begin{proposition}[Necessary dependence on a vanishing positivity floor]\label{prop:q-sharp}
For the binary audit $Z=y+tE$, suppose a pre-report preparation can have sign probabilities $(1-q,q)$ and a weak sign report has likelihoods $1-\gamma,\gamma$, with $0<\gamma<1/2$. On a positive report the smaller posterior probability is
\[
 a_q=\frac{q\gamma}{(1-q)(1-\gamma)+q\gamma}.
\]
Its residual variance is $4a_q(1-a_q)t^2=O(qt^2)$. If the admitted family also contains the sign-reversed preparations, its conditional minimax radius is of order $t^2$ for small $q$. Consequently the best uniform ratio in Theorem~\ref{thm:fiber} is at least of order $q^{-1}$. The loss of uniformity as $q\downarrow0$ in this mark mechanism is necessary.
\end{proposition}
\begin{proof}
Bayes' rule gives the displayed fraction. The two extreme preparations, conditional on the same weak positive report, produce audit means tending to $y+t$ and $y-t$. Any center has squared distance at least a fixed multiple of $t^2$ from one of them. The least variance is bounded by the first preparation's variance. Their ratio has the stated order. The same calculation is obtained from diagonal qubit states and a weak $\sigma_z$ effect. It is a subfamily of an instrument class that may also contain noncommuting commands. This proves a limitation of uniform proxy calibration, not a complete minimax law in the joint $q,N,M$ limit.
\end{proof}
